% Part: computability
% Chapter: recursive-functions
% Section: primitive-recursion

\documentclass[../../../include/open-logic-section]{subfiles}

\begin{document}

\olfileid[tr]{cmp}{rec}{pre}
\olsection{İlkel Özyineleme}

Doğal sayıların bir özelliği, her doğal sayıya $0$'dan başlayıp ardıl
işlemi~$+1$ sonlu sayıda uygulanarak ulaşılabilmesidir: Her doğal sayı
ya $0$'dır ya da $0$'ın ardılının ... ardılıdır. Bu olgudan
yararlanarak bir $h\colon\Nat\to\Nat$ fonksiyonu belirtmenin bir yolu
şudur: (a) $0$ girdisi için $h$'nin değerini belirtmek ve (b) $h(x)$
değeri verildiğinde $h(x+1)$ değerinin nasıl hesaplanacağını belirtmek.
(a), $h(0)$ değerini doğrudan verir; dolayısıyla $h$, $0$ için
tanımlıdır. Şimdi (b)'deki yönergeyi $x=0$ için kullanarak
$h(1)=h(0+1)$ değerini $h(0)$'dan hesaplarız. Aynı yönergeyi $x=1$
için kullanıp $h(2)=h(1+1)$ değerini $h(1)$'den hesaplarız ve böyle
sürdürürüz. Her doğal sayı~$x$ için sonunda $h(x+1)$'i $h(x)$'den
tanımladığımız adıma ulaşırız; böylece $h$, bütün $x\in\Nat$ için
tanımlanır.

Örneğin $h\colon\Nat\to\Nat$ fonksiyonunu şu iki denklemle
belirttiğimizi varsayalım:
\begin{align*}
h(0) & =  1\\
h(x+1) & =  2 \cdot h(x).
\end{align*}
Çarpmayı zaten biliyorsak bu denklemler yukarıdaki (a) ve (b) için
gereken bilgiyi verir. İkinci denklemi art arda uygulayarak
\begin{align*}
  h(1) & = 2\cdot h(0) = 2,\\
  h(2) & = 2\cdot h(1) = 2\cdot 2,\\
  h(3) & = 2 \cdot h(2) = 2\cdot 2 \cdot 2,\\
  & \vdots
\end{align*}
elde ederiz. Belirttiğimiz fonksiyonun $h(x)=2^x$ olduğunu görürüz.

Doğal sayıların bu ayırt edici özelliği, bu iki ölçütü karşılayan
yalnızca bir $h$ fonksiyonu bulunduğunu güvence altına alır. Böyle bir
denklem çiftine $h$ fonksiyonunun \emph{ilkel özyinelemeyle tanımı}
denir. ``Özyinelemeli'' denmesinin nedeni tanımın, özellikle ikinci
denklemin, sağ tarafında $h$'nin kendisini içermesidir. ``İlkel''
denmesinin nedeni ise $h(x+1)$ tanımlanırken yalnızca bir önceki
değer~$h(x)$'in kullanılmasıdır. Bu, $\Nat$ üzerinde bir fonksiyonu
özyinelemeyle tanımlamanın en basit yoludur.

Toplama ve çarpma gibi daha temel fonksiyonları da ilkel özyinelemeyle
tanımlayabiliriz. Ancak bu durumda söz konusu fonksiyonlar iki
yerlidir. Girdi yerlerinden birini sabitler, özyineleme için ötekini
kullanırız. Örneğin $\Add(x,y)$ tanımlamak için $x$'i sabitleyip önce
$y=0$ değerini, ardından $y+1$ değerini $y$ cinsinden tanımlarız.
$x$ sabit olduğu için tanımlayıcı denklemlerin her iki tarafında da
görünür:
\begin{align*}
\Add(x,0) & = x\\
\Add(x,y+1) & = \Add(x,y)+1.
\end{align*}
Bu denklemler $\Add$ değerini bütün $x$ ve $y$ için belirtir. Örneğin
$\Add(2,3)$ bulmak için $x=2$ ile tanımlayıcı denklemleri uygularız:
Birinciyle $\Add(2,0)=2$; ikinciyi art arda kullanarak
$\Add(2,1)=2+1=3$, $\Add(2,2)=3+1=4$ ve $\Add(2,3)=4+1=5$ buluruz.

$\Add$ tanımının ikinci denkleminin sağ tarafında $+$ kullandık; fakat
yalnızca $1$ eklemek için. Başka bir deyişle ardıl fonksiyonu
$\Succ(z)=z+1$ kullandık ve $\Add(x,y+1)$ değerini tanımlamak için onu
önceki $\Add(x,y)$ değerine uyguladık. Bu nedenle özyinelemeli tanımı,
önceki değere uyguladığımız tek bir fonksiyon cinsinden verilmiş
düşünebiliriz. Ancak fonksiyonun yalnızca önceki değere değil, $x$ ve
$y$'ye de bağlı olmasına izin vermek sakıncalı değildir, hatta bazen
gereklidir. Şunu ele alalım:
\begin{align*}
  \Mult(x,0) & = 0 \\
  \Mult(x,y+1) & = \Add(\Mult(x,y),x).
\end{align*}
Bu, $\Add$ fonksiyonunu hem önceki değer $\Mult(x,y)$'ye hem ilk
girdi~$x$'e uygulayarak verilen $\Mult$ fonksiyonunun ilkel
özyinelemeli tanımıdır. $\Mult(x,y)$ değerini bütün $x$ ve $y$ için
tanımlar. Örneğin $\Mult(2,3)$, sırasıyla $\Mult(2,0)$,
$\Mult(2,1)$, $\Mult(2,2)$ ve $\Mult(2,3)$ hesaplanarak belirlenir:
\begin{align*}
  \Mult(2,0) & = 0\\
  \Mult(2,1) & = \Mult(2,0+1) =
  \Add(\Mult(2,0), 2) = \Add(0, 2) = 2\\
  \Mult(2,2) & = \Mult(2,1+1) =
  \Add(\Mult(2,1), 2) = \Add(2, 2) = 4\\
  \Mult(2,3) & = \Mult(2,2+1) =
  \Add(\Mult(2,2), 2) = \Add(4, 2) = 6.
\end{align*}

Genel örüntü şudur: Bir $h(x_0,\dots,x_{k-1},y)$ fonksiyonunun ilkel
özyinelemeli tanımını vermek için iki denklem veririz. Birincisi
$h(x_0,\dots,x_{k-1},0)$ değerini $h$'ye başvurmadan tanımlar.
İkincisi $h(x_0,\dots,x_{k-1},y+1)$ değerini
$h(x_0,\dots,x_{k-1},y)$, öteki girdiler $x_0,\dots,x_{k-1}$ ve $y$
cinsinden tanımlar. İkinci denklemde yalnızca $h$'nin hemen önceki
değeri kullanılabilir. Bu iki denklemin sağ taraflarında verilen
işlemleri kendileri $f$ ve $g$ fonksiyonları olarak düşünürsek,
$h$'yi ilkel özyinelemeyle tanımlamanın genel örüntüsü şöyledir:
\begin{align*}
  h(x_0, \dots, x_{k-1}, 0) & = f(x_0, \dots, x_{k-1})\\
  h(x_0, \dots, x_{k-1}, y+1) & =
  g(x_0, \dots, x_{k-1}, y, h(x_0, \dots, x_{k-1}, y)).
\end{align*}
$\Add$ durumunda $k=1$, $f(x_0)=x_0$ (özdeşlik fonksiyonu) ve
$g(x_0,y,z)=z+1$ (üçüncü girdisinin ardılını döndüren üç yerli
fonksiyon) olur:
\begin{align*}
  \Add(x_0, 0) & = f(x_0) = x_0\\
  \Add(x_0, y+1) & = g(x_0, y, \Add(x_0, y)) =
  \Succ(\Add(x_0, y)).
\end{align*}
$\Mult$ durumunda $f(x_0)=0$ (her zaman $0$ döndüren sabit fonksiyon)
ve $g(x_0,y,z)=\Add(z,x_0)$ (son ve ilk girdisinin toplamını döndüren
üç yerli fonksiyon) olur:
\begin{align*}
  \Mult(x_0,0) & = f(x_0) = 0 \\
  \Mult(x_0,y+1) & = g(x_0, y, \Mult(x_0,y)) =
  \Add(\Mult(x_0,y), x_0).
\end{align*}

\end{document}
